%% Report template (REVTeX 4.2). generated/ and figures/ are written by ../analysis/run_all.py.
\documentclass[aps,pra,reprint,superscriptaddress,nofootinbib,floatfix]{revtex4-2}

\usepackage{graphicx}
\usepackage{amsmath,amssymb}
\usepackage{bm}
\usepackage{xcolor}
\usepackage[colorlinks=true,linkcolor=blue!50!black,citecolor=blue!50!black,urlcolor=blue!50!black]{hyperref}

% generated by analysis/report.py -- do not edit
\newcommand{\calA}{-12.34}
\newcommand{\calB}{0.24749}
\newcommand{\calChi}{4.2}
\newcommand{\calNdf}{2}
\newcommand{\resFWHMsix}{6.9}
\newcommand{\resFWHMthree}{8.3}
\newcommand{\gainShiftBa}{0.004}
\newcommand{\gainShiftNa}{0.16}
\newcommand{\gainShiftCs}{2.0}
\newcommand{\gainDriftA}{0.3}
\newcommand{\gainDriftB}{1.0}
\newcommand{\beamStability}{0.7}
\newcommand{\devEightyOne}{-0.2}
\newcommand{\Szero}{3.24}
\newcommand{\dSzero}{0.09}
\newcommand{\psiZero}{5.12 \pm 0.08}
\newcommand{\wEdge}{0.80 \pm 0.05}
\newcommand{\haloAmp}{0.104 \pm 0.024}
\newcommand{\haloLam}{2.97 \pm 0.27}
\newcommand{\leakRate}{0.42 \pm 0.10}
\newcommand{\beamChi}{6.0}
\newcommand{\beamNdf}{4}
\newcommand{\pileup}{0.69}
\newcommand{\dpileup}{0.34}
\newcommand{\rateZero}{950.5}
\newcommand{\removal}{0.423}
\newcommand{\removalErr}{0.009}
\newcommand{\odTrans}{2.14}
\newcommand{\odPhoto}{1.30}
\newcommand{\dodPhoto}{0.30}
\newcommand{\yOffset}{-0.4}
\newcommand{\dyOffset}{0.2}
\newcommand{\removalZeroDef}{0.435}
\newcommand{\idod}{0.40}
\newcommand{\didod}{0.05}
\newcommand{\shadowMeas}{+0.67}
\newcommand{\shadowErr}{0.50}
\newcommand{\shadowPredCentred}{-7.0}
\newcommand{\shadowPredNominal}{+0.6}
\newcommand{\alphaMean}{0.78}
\newcommand{\alphaStat}{0.05}
\newcommand{\alphaSyst}{0.11}
\newcommand{\alphaTot}{0.12}
\newcommand{\alphaConstChi}{9.5}
\newcommand{\alphaMeanAlt}{1.35}
\newcommand{\alphaStatAlt}{0.09}
\newcommand{\alphaSystAlt}{0.49}
\newcommand{\knShapeTotChi}{7.5}
\newcommand{\knShapeTotP}{0.18}
\newcommand{\thShapeTotChi}{29.8}
\newcommand{\thShapeTotP}{1.7\times10^{-5}}
\newcommand{\reVal}{2.49}
\newcommand{\reStat}{0.08}
\newcommand{\reSyst}{0.17}
\newcommand{\knShapeChi}{9.5}
\newcommand{\knShapeNdf}{5}
\newcommand{\knShapeP}{0.09}
\newcommand{\thShapeChi}{39.1}
\newcommand{\thShapeNdf}{5}
\newcommand{\thShapeP}{2.2\times10^{-7}}
\newcommand{\knShapeC}{0.78}
\newcommand{\knShapeDC}{0.05}
\newcommand{\thShapeC}{0.33}
\newcommand{\thShapeDC}{0.02}
\newcommand{\knAbsChi}{7.5}
\newcommand{\knAbsNdf}{6}
\newcommand{\knAbsP}{0.28}
\newcommand{\thAbsChi}{24.7}
\newcommand{\thAbsNdf}{6}
\newcommand{\thAbsP}{3.9\times10^{-4}}
\newcommand{\thShapeChiAlt}{34.3}
\newcommand{\knShapeChiAlt}{8.8}
\newcommand{\mFix}{522}
\newcommand{\mFixStat}{7}
\newcommand{\mFixSyst}{16}
\newcommand{\mFixChi}{4.6}
\newcommand{\mFixNdf}{5}
\newcommand{\mFree}{504}
\newcommand{\dmFree}{19}
\newcommand{\EzeroFree}{673}
\newcommand{\dEzeroFree}{13}
\newcommand{\mAll}{512}
\newcommand{\dmAll}{15}
\newcommand{\noShiftChi}{4252}
\newcommand{\noShiftNdf}{6}
\newcommand{\mLinT}{511}
\newcommand{\dmLinT}{18}
\newcommand{\mLinG}{513}
\newcommand{\dmLinG}{14}
\newcommand{\EzeroLinT}{671}
\newcommand{\dEzeroLinT}{12}
\newcommand{\nFit}{6}
\newcommand{\selfAbsMin}{0.38}
\newcommand{\selfAbsMax}{0.41}
\newcommand{\msMin}{0.30}
\newcommand{\msMax}{0.56}
\newcommand{\thetaRmsMin}{4.0}
\newcommand{\thetaRmsMax}{4.7}
\newcommand{\SBmin}{0.08}
\newcommand{\SBmax}{0.35}
\newcommand{\epsFEPsix}{0.296}
\newcommand{\epsFEPnin}{0.689}
\newcommand{\PTsix}{0.47}
\newcommand{\PTnin}{0.83}
\newcommand{\epsDirect}{0.325}
\newcommand{\odKN}{1.67}
\newcommand{\matSteelT}{0.78}
\newcommand{\matAlT}{0.56}
\newcommand{\matCuT}{0.82}
\newcommand{\matAlG}{2.61}
\newcommand{\odAlT}{4.04}
\newcommand{\odCuT}{2.02}
\newcommand{\KNtotal}{0.2562}
\newcommand{\Lone}{28.5}
\newcommand{\dLone}{0.5}
\newcommand{\Ltwo}{48.5}
\newcommand{\dLtwo}{0.5}
\newcommand{\rdist}{20.0}
\newcommand{\aperture}{2.54}
\newcommand{\daperture}{0.30}
\newcommand{\dOmega}{0.051}
\newcommand{\angleUnc}{0.6}
\newcommand{\faceUnc}{1.5}
\newcommand{\relEffSyst}{5}
\newcommand{\removalSystRel}{2}
\newcommand{\histories}{300{,}000}
\newcommand{\mFixRel}{3.3}
\newcommand{\comptonWavelength}{2.4263}
\newcommand{\reCODATA}{2.818}
\newcommand{\yKNmin}{713}
\newcommand{\yKNmax}{1498}
\newcommand{\nGoodAngles}{6}
\newcommand{\rmsCalResid}{0.52}
\newcommand{\alphaSigmaKN}{1.9}
\newcommand{\thShapeZ}{4.3}
\newcommand{\knShapeTotNdf}{5}
\newcommand{\fitChiMin}{0.8}
\newcommand{\fitChiMax}{1.2}
\newcommand{\alphaMinFit}{0.39}
\newcommand{\alphaMaxFit}{0.93}
\newcommand{\alphaNinety}{0.39}
\newcommand{\yNinety}{278}
\newcommand{\dyNinety}{111}
\newcommand{\bkgNinety}{3324}
\newcommand{\NTmin}{0.74}
\newcommand{\NTmax}{0.93}


\newcommand{\Cs}{$^{137}$Cs}
\newcommand{\Ba}{$^{133}$Ba}
\newcommand{\Na}{$^{22}$Na}
\newcommand{\NaI}{NaI(Tl)}
\newcommand{\dsdo}{d\sigma/d\Omega}
\newcommand{\mec}{m_ec^2}

\begin{document}

\title{Compton scattering of 662-keV $\gamma$ rays from a steel target:\\
energy--angle relation, electron rest energy, and differential cross section}

\author{Author One}
\author{Author Two}
\author{Author Three}
\affiliation{Department of Physics, [Institution], [City, Country]}
\date{\today}

\begin{abstract}
Abstract
\end{abstract}

\maketitle

\section{Introduction}
\label{sec:intro}

Introduction

\section{Theory and expected results}
\label{sec:theory}

\subsection{Kinematics of Compton scattering}

%------------------------------------------------

\subsection{Thomson and Klein--Nishina cross sections}

%------------------------------------------------

\subsection{Interaction of the scattered photons with the target and the detector}

%------------------------------------------------

\subsection{From counts to a cross section}
\label{sec:measeq}

%------------------------------------------------

\subsection{Expected results}

%------------------------------------------------

\section{Experimental setup and data}
\label{sec:setup}

\subsection{Apparatus}

%------------------------------------------------

\subsection{Data sets}

%------------------------------------------------

\section{Analysis}
\label{sec:analysis}

The analysis is implemented in Python. A single configuration file holds every input and its uncertainty, and one script regenerates all numbers, tables and figures of this report.

\subsection{Energy calibration and resolution}
\label{sec:calib}

%------------------------------------------------

\subsection{The measured spectra}
\label{sec:spectra}

%------------------------------------------------

\subsection{The primary beam}
\label{sec:beam}

%------------------------------------------------

\subsection{Target thickness}
\label{sec:target}

%------------------------------------------------

\subsection{Detector efficiency}
\label{sec:eff}

%------------------------------------------------

\subsection{Forward-folding model of the scattered spectrum}
\label{sec:fold}

%------------------------------------------------

\subsection{Signal extraction}
\label{sec:extract}

%------------------------------------------------

\section{Uncertainties and errors}
\label{sec:uncertainties}

%------------------------------------------------

\subsection{Statistical uncertainties}

%------------------------------------------------

\subsection{Signal-extraction uncertainty}

%------------------------------------------------

\subsection{Systematic uncertainties}

%------------------------------------------------

\subsection{Correlations and combination}

%------------------------------------------------

\subsection{Energy uncertainties}

%------------------------------------------------

\subsection{Errors identified and corrected}
\label{sec:errors}

%------------------------------------------------

\subsection{Compton kinematics and the electron rest energy}
\label{sec:kin}

%------------------------------------------------

\subsection{Differential cross section}

%------------------------------------------------

\subsection{Dependence on the target assumptions}
\label{sec:norm}

%------------------------------------------------

\section{Discussion}
\label{sec:discussion}

%------------------------------------------------

\section{Conclusion}
\label{sec:conclusion}

Conclusion

\appendix

\section{Calibration cross-check}
\label{app:calcheck}

%------------------------------------------------

\section{Reproducing the results}
\label{app:code}

%------------------------------------------------

\bibliography{references}

\end{document}
